\documentclass{article}
\usepackage[utf8]{inputenc}
\usepackage[T1]{fontenc}
\usepackage[margin=2cm]{geometry}
\usepackage{hyperref}
\sloppy
\begin{document}
\section*{Historical source}
Encyclopaedia Britannica scan, source row 1473. Parse 3.5 extracted the text. Super 3.5 VL corrected it from an enlarged crop of each block. It can still contain recognition errors.
Page headers, footers, and page numbers are removed.

\subsection*{ESARHADDON (E0123)}
ESARHADDON, the son of Sennacherib, and his
successor in the kingdom of Assyria. He is said to
have reigned 29 years at Nineveh, from the year 3294
to 3322; besides which he reigned 13 years at Baby-
lon, in all 42 years. He died in the year of the world
3336, and was succeeded by Saosduchinus. Esar-
haddon, in the opinion of Sir Isaac Newton, seems to
be the Sardanapalus who died, as Ctesias says, of
old age, after the revolt of Syria; the name Sardana-
palus being derived from Asarhaddon Pul.

\subsection*{ESCALADE (E0124)}
ESCALADE, or SCALADE, a furious attack of a
wall or a rampart; carried on with ladders, to pass the
ditch or mount the rampart; without proceeding in
form, breaking ground, or carrying on regular works
to secure the men.

When the troops are prepared to pass the ditch,
either with the assistance of boards, hurdles, and fascines, when it is muddy, or with small boats of tin, or
baskets covered with skins or oil-cloth, when it is deep
and filled with water, a party must be placed on the
counterscarp, opposite to the landing-place, ready to
fire at the garrison if they are alarmed, and oppose the
mounting on the rampart. If the ditch is dry, the
ladders are fixed in some place farther distant from the
centre; and as soon as they get upon the rampart,
they put themselves in order to receive the enemy; if
the centre should be furbished and likely overcome,
the detachment hastens to break open the gate, and to
let in the rest of the party. If the ditch is wet, the
rampart high, and provided with a revetement, it will
be difficult to surprise the town in this way; but if
there is no revetement, the troops may hide themselves
along the outside of the rampart till all are over. Since the invention and use of gunpowder, and the
walls of cities have been flanked, they are seldom taken
by escalade.

\subsection*{ESCALDOLIN (E0125)}
ESCALDOLIN, In January, a grass off the man
magnetor, belonging to the presumiric data of
plants. The fruit is bluebar and poly/permons; the
petal diftants and longest-lived; the figures headed.

\subsection*{ESCAPE (E0126)}
ESCAPE, in law, a violent or privy evasion out of
some lawful restraint, without being delivered by due
course of law. There are two sorts of escapes, volun-
tary and negligent. Voluntary, when a man arrests
another for felony, or other crime, and afterwards lets
him go freely by consent; in which case, the party
that permits such escape is held guilty, committed,
and must answer for it. Negligent escape, on the con-
trary, is where one is arrested, and afterwards escapes
against the will of the person that arrested him, and is
not pursued with fresh suit, and retaken before the
person purloining hath lost fight of him. By stat. 8 and
9 Will. III. c. 26. the keepers of prisons conniving at
escapes shall forfeit 500l.; and in civil cases the sher-
iff is answerable for the debt.

\subsection*{ESCHALOT (E0127)}
ESCHALOT, or Shallot. See Allium.

\subsection*{ESCHAR (E0128)}
ESCHAR, in surgery, the crust or scab occasioned
by burns or caustic medicines.

\subsection*{ESCHARA (E0129)}
ESCHARA, in natural history, the name of a species of coralline, \&c. the characters of which are these: they are of a stony or coral-like hardness, and resemble a woven cloth in their texture; and the microscope informs us, that they consist of arrangements of very small cells, whose surfaces appear much in that form.

Linnaeus makes it a species of millepora, in the class of lithophytes. See Plate CXLVIII. fig. 9. to.

The narrow-leaved hornwrack, fig. 9, divides as it
rises, into narrow leaves made up of regular rows of
oblong square-shaped cells placed alternately by each
other, and opposite to an equal number on the other
side of the leaf, like an honeycomb; from these leaves
proceed other still smaller foliaceous ramifications,
many of which seem to be connected at the lower part
by tubuli, as in the corallines; by which means they
can ply to and fro more freely in the waters.—c. Gives
the natural appearance of this coralline. E. represents
two leaves with their tubuli and cells magnified. E 1
is a cross section of one of the leaves at E, flowing the
partition and inner form of their cells.

The broad-leaved hornwrack, fig. to. when fresh taken out of the tea, is of a spongy soft texture, and smells very filthy; but when it has lain for some time on the floor, it becomes stiff and horny, like some fort of withered leaves. Both surfaces, when examined by glasses, appear to be covered with cells; and, when a piece of it is cut across, one may discover the thin membrane that serves as a base to the cells of each fur- face. The form of the cells is very remarkable, each

The form of the cells is very remarkable, each one being arched at the top, and contracted a little at the lower part of the sides to make way for the arches of the two next adjoining cells; so that by this particular contruction no room is lost. The entrance of the cells is immediately under the arch of each cell, and the walls of the cells seem to be fortified with spines. Jussieu discovered small polypes extending themselves out of these cells, which he has described in the memoirs of the Academy of Sciences 1742.—f, Gives the natural appearance of a leafy branch of this coralline. F is a part of a leaf magnified to show the superficial figure of the cells, and the manner in which they are disposed. F f shows a cross section of a leaf, and discovers the several partitions of the cells.

At the entrance of many of the cells a small testaceous body like a bivalve shell is discovered : F 2, the figure of the cell, with the shell in it ; it is of a transparent amber colour, so clear that one may fee the dead animal through it, represented by the black spot.

\subsection*{ESCHEAT (E0130)}
ESCHEAT, in law, signifies any lands or tenements that casually fall to a lord within his manor. It is one of the consequences of tenure in chivalry. (See Feodal System, Knight-Service, and Tenure). It is the determination of the tenure or dissolution of the mutual bond between the lord and tenant, from the extinction of the blood of the latter by either natural or civil means: if he died without heirs of his blood, or if his blood was corrupted and tainted by commission of treason or felony; whereby every inheritable quality was entirely blotted out and abolished. In such cases the land escheated or fell back to the lord of the fee; that is, the tenure was determined by breach of the original condition, expressed or implied in the feudal donation. In the one case, there were no heirs substituting of the blood of the first feudatory or purchaser, to which heirs alone the grant of the feud extended: in the other, the tenant, by perpetrating an atrocious crime, showed that he was no


\section*{Source record}
\begin{description}
\item[Scan ID] \texttt{britannica:309f765b1ecc83e0a7f7762e44b4a927a7555443:train:1473}
\item[Image file] \texttt{page-1473-b3b602ef6a3d.jpg}
\item[Image SHA-256] \texttt{b3b602ef6a3d4961fd1926a15d634a40ec929c15b00e2f7cc5109b3f7180daac}
\item[Dataset] NationalLibraryOfScotland/encyclopaedia\_britannica\_illustrated at revision \texttt{309f765b1ecc83e0a7f7762e44b4a927a7555443} (CC0-1.0)
\item[Extraction] model \texttt{nvidia/nemotron-3.5-parse}, raw reply \texttt{extract/raw/page-1473-b3b602ef6a3d.json}
\item[Raw reply SHA-256] \texttt{15eb042f132389fdb22ccc015e1640219500d9804eae7e0608b3fc40ac6c6c1a}
\item[Correction] model \texttt{nvidia/Nemotron-3.5-Super-VL}, block crops, raw replies \texttt{correct/raw/crop\_nothink/page-1473-b3b602ef6a3d-b*.json}; parse\_fallback: 1, super\_vl: 13
\item[Running header] ESC 1723 I ESC
\end{description}
\end{document}
